% predicate-logic-inference.tex

%%%%%%%%%%%%%%%%%%%%
\begin{frame}{}
  \begin{center}
    {\Large ``再回到语法''}

    \vspace{0.30cm}
    \fig{width = 0.60\textwidth}{figs/syntax-semantics}
    \vspace{0.30cm}

    \pause
    {\large 暂时\cyan{\bf 忘掉语义}!!! 一起来研究公式之间的\blue{\bf 推导}关系}
  \end{center}
\end{frame}
%%%%%%%%%%%%%%%%%%%%

%%%%%%%%%%%%%%%%%%%%
\begin{frame}{}
  \begin{center}
    {\Large 一阶谓词逻辑的自然推理 (演绎; 推演) 系统} \\[10pt]
  \end{center}
\end{frame}
%%%%%%%%%%%%%%%%%%%%

%%%%%%%%%%%%%%%%%%%%
\begin{frame}{\texorpdfstring{\forallxelim}{forall x-elim}}
  \begin{center}
    \[
      \nd{\forall x\; \phi}{\phi[t/x]}{\forallxelim}
    \]

    where $t$ is \red{free} for $x$ in $\phi$
  \end{center}

  \pause
  \[
    \forall x\; \blue{P(x)} \vdash P(\blue{c}) \quad (c \;\text{是任意常元符号})
  \]
  \pause
  \[
    \forall x\; \blue{\exists y\; (x < y)} \vdash
      \exists y\; (\blue{c} < y) \quad (c \;\text{是任意常元符号})
  \]
  \pause
  \[
    \forall x\; \blue{\exists y\; (x < y)} \vdash
      \exists y\; (\blue{z} < y)
      \quad (z \neq y \;\text{是任意变元符号})
  \]
  \pause
  \[
    \forall x\; \blue{\exists y\; (x < y)} \nvdash
      \exists y.\; (\red{y} < y)
      \quad (y \;\text{is {\it not} free for } x \text{ in } \phi)
  \]
\end{frame}
%%%%%%%%%%%%%%%%%%%%

%%%%%%%%%%%%%%%%%%%%
\begin{frame}{\texorpdfstring{\forallxelim}{forall x-elim} 推理规则的应用}
  \[
    \ndnoname{\text{每个人都是要死的 \qquad 苏格拉底是人}}{\text{苏格拉底是要死的}}
  \]

  \vspace{0.80cm}
  \[
    \ndnoname{\forall x\; (H(x) \to M(x)) \qquad H(s)}{M(s)}
  \]
\end{frame}
%%%%%%%%%%%%%%%%%%%%

%%%%%%%%%%%%%%%%%%%%
\begin{frame}{\texorpdfstring{\forallintro}{forall-intro}}
  \[
    \nd{
      \fbox{%
        $\begin{array}{c}
          x_{0} \;\text{\red{fresh}} \\[1mm]
          \vdots \\[1mm]
          \phi[x_0/x]
        \end{array}$%
      }
    }{\forall x\;\phi}{\forallxintro}
  \]

  \vspace{2em}
  \begin{center}
    $x_0$ is a \red{\bf fresh} variable,
    representing an \blue{\bf arbitrary} term
  \end{center}
\end{frame}
%%%%%%%%%%%%%%%%%%%%

%%%%%%%%%%%%%%%%%%%%
\begin{frame}{}
  \begin{center}
    Everyone loves Bob Dylan. \\[8pt]
    Therefore, everyone loves themselves.
  \end{center}

  \pause
  \[
    \forall x\; L(x, B) \vdash \forall x\; L(x, x)
  \]

  \vspace{2em}
  \begin{logicproof}{2}
    \uncover<3->{\forall x\; L(x, B)} & \uncover<3->{premise} \\
    \uncover<5->{
    \begin{subproof}
      \uncover<6->{B} & \uncover<6->{\red{fresh}} \\
      \uncover<7->{L(B, B)} & \uncover<7->{\forallxelim 1}
    \end{subproof}}
    \uncover<4->{\forall x\; L(x, x)} & \uncover<8->{\forallxintro 2--3}
  \end{logicproof}
\end{frame}
%%%%%%%%%%%%%%%%%%%%

%%%%%%%%%%%%%%%%%%%%
\begin{frame}{}
  \begin{exampleblock}{\texorpdfstring{$\forallintro$}{forall-intro} 推理规则的应用}
    \[
      \forall x(P(x)\to Q(x)),\ \forall x P(x) \vdash \forall x Q(x)
    \]
  \end{exampleblock}

  \pause
  \vspace{2em}
  \begin{logicproof}{2}
    \uncover<2->{\forall x(P(x)\to Q(x))} & \uncover<2->{premise} \\
    \uncover<2->{\forall x P(x)} & \uncover<2->{premise} \\
    \uncover<4->{
    \begin{subproof}
      \uncover<5->{x_0} & \uncover<5->{\blue{\text{fresh}}} \\
      \uncover<7->{P(x_0)\to Q(x_0)} & \uncover<7->{\forallxelim 1} \\
      \uncover<8->{P(x_0)} & \uncover<8->{\forallxelim 2} \\
      \uncover<6->{Q(x_0)} & \uncover<9->{\toelim 4, 5}
    \end{subproof}}
    \uncover<3->{\forall x Q(x)} & \uncover<10->{\forallxintro 3--6}
  \end{logicproof}
\end{frame}
%%%%%%%%%%%%%%%%%%%%

%%%%%%%%%%%%%%%%%%%%
\begin{frame}{}
  % P111 in LICS
  \begin{exampleblock}{\texorpdfstring{$\forallintro$}{forall-intro} 推理规则的应用}
    \[
      \Big\{P(t), \forall x\; (P(x) \to \lnot Q(x))\Big\}
        \vdash \lnot Q(t)
    \]
  \end{exampleblock}

  \pause
  \vspace{2em}
  \begin{logicproof}{2}
    P(t) & premise \\
    \forall x(P(x)\to \lnot Q(x)) & premise \\
    \uncover<4->{P(t)\to \lnot Q(t)} & \uncover<4->{\forallxelim 2} \\
    \uncover<3->{\lnot Q(t)} & \uncover<5->{\toelim 3, 1}
  \end{logicproof}
\end{frame}
%%%%%%%%%%%%%%%%%%%%

%%%%%%%%%%%%%%%%%%%%
\begin{frame}{\texorpdfstring{\existsintro}{exists-intro}}
  \begin{center}
    \[
      \nd{\phi[t/x]}{\exists x\; \phi}
        {\existsxintro}
    \]
  \end{center}

  \pause
  \[
    P(\blue{c}) \vdash \exists x\; P(x)
      \quad c \text{ 是任意常元符号}
  \]

  \pause
  \[
    \teal{\forall y\; (y = y)} \nvdash
      \exists x\; \teal{\forall y\; (x = y)}
  \]
\end{frame}
%%%%%%%%%%%%%%%%%%%%

%%%%%%%%%%%%%%%%%%%%
\begin{frame}{\texorpdfstring{$\exists\text{-elim}$}{exists-elim}}
  \[
    \nd{
      \begin{array}{c}
        \exists x\,\phi
        \qquad
        \fbox{%
          $\begin{array}{c}
            x_0 \ \text{fresh} \\[1mm]
            \phi[x_0/x] \\[1mm]
            \vdots \\[1mm]
            \phi
          \end{array}$%
        }
      \end{array}
    }{\phi}{\existsxelim}
  \]

  \vspace{1em}
  \begin{center}
    \red{$x_{0}$ {\it cannot} occur outside its box
      (in particular, cannot occur in $\phi$)}
  \end{center}
\end{frame}
%%%%%%%%%%%%%%%%%%%%

%%%%%%%%%%%%%%%%%%%%
\begin{frame}{}
  \begin{center}
    Bob Dylan is a poet. \\[5pt]
    Someone is not a poet. \\[5pt]
    Therefore, Bob Dylan is both a poet and not a poet.
  \end{center}

  \pause
  \[
    \left\{P(B), \exists x\; \lnot P(x) \vdash
      P(B)\land \lnot P(B) \right\}
  \]

  \vspace{2em}
  \begin{logicproof}{2}
    \uncover<3->{P(B)} & \uncover<3->{premise} \\
    \uncover<3->{\exists x \lnot P(x)} & \uncover<3->{premise} \\
    \uncover<5->{
    \begin{subproof}
      \uncover<6->{B\; (\text{\blue{fresh}}),\; \lnot P(B)} & \uncover<6->{assumption} \\
      \uncover<7->{P(B) \land \lnot P(B)} & \uncover<7->{\landintro 1, 3}
    \end{subproof}}
    \uncover<4->{P(B) \land \lnot P(B)} & \uncover<8->{\existsxelim 2, 3--4}
  \end{logicproof}
\end{frame}
%%%%%%%%%%%%%%%%%%%%

%%%%%%%%%%%%%%%%%%%%
\begin{frame}{}
  \begin{exampleblock}{$\existsintro$ 推理规则的应用}
    \[
      \forall x\; P(x) \vdash \exists x\; P(x)
    \]
  \end{exampleblock}

  \vspace{2em}
  \begin{logicproof}{1}
    \uncover<2->{\forall x \phi} & \uncover<2->{premise} \\
    \uncover<4->{\phi[x/x]} & \uncover<4->{\forallxelim 1} \\
    \uncover<3->{\exists x \phi} & \uncover<5->{\existsxintro 2}
  \end{logicproof}
\end{frame}
%%%%%%%%%%%%%%%%%%%%

%%%%%%%%%%%%%%%%%%%%
\begin{frame}{}
  \begin{exampleblock}{$\forall, \exists$\text{-推理规则的应用}}
    \[
      \Big\{\forall x\; (P(x) \to Q(x)), \exists x\; P(x)\Big\}
        \vdash \exists x\; Q(x)
    \]
  \end{exampleblock}

  \vspace{2em}
  \begin{logicproof}{2}
    \uncover<2->{\forall x(P(x)\to Q(x))} & \uncover<2->{premise} \\
    \uncover<3->{\exists x P(x)} & \uncover<3->{premise} \\
    \uncover<5->{
    \begin{subproof}
      \uncover<6->{x_0\; (\text{\blue{fresh}}),\; P(x_0)} & \uncover<6->{assumption} \\
      \uncover<8->{P(x_0)\to Q(x_0)} & \uncover<8->{\forallxelim 1} \\
      \uncover<9->{Q(x_0)} & \uncover<9->{\toelim 4, 3} \\
      \uncover<7->{\exists x Q(x)} & \uncover<10->{\existsxintro 5}
    \end{subproof}}
    \uncover<4->{\exists x Q(x)} & \uncover<11->{\existsxelim 2, 3--6}
  \end{logicproof}
\end{frame}
%%%%%%%%%%%%%%%%%%%%

%%%%%%%%%%%%%%%%%%%%
\begin{frame}{}
  \begin{exampleblock}{$\forall, \exists$\text{-推理规则的 (\red{\bf 错误}) 应用}}
    \[
      \Big\{\forall x\; (P(x) \to Q(x)), \exists x\; P(x)\Big\}
        \vdash \exists x\; Q(x)
    \]
  \end{exampleblock}

  \vspace{2em}
  \begin{logicproof}{2}
    \uncover<2->{\forall x\; (P(x) \to Q(x))} & \uncover<2->{premise} \\
    \uncover<3->{\exists x\; P(x)} & \uncover<3->{premise} \\
    \uncover<5->{
    \begin{subproof}
      \uncover<6->{x_0\; (\text{\blue{fresh}}),\; P(x_0)} & \uncover<6->{assumption} \\
      \uncover<7->{P(x_0) \to Q(x_0)} & \uncover<7->{\forallxelim 1} \\
      \uncover<8->{Q(x_0)} & \uncover<8->{\toelim 4, 3}
    \end{subproof}}
    \uncover<9->{Q(x_0)} & \uncover<9->{\existsxelim 2, 3--5} \\
    \uncover<4->{\exists x\; Q(x)} & \uncover<10->{\existsxintro 6}
  \end{logicproof}
\end{frame}
%%%%%%%%%%%%%%%%%%%%

%%%%%%%%%%%%%%%%%%%%
\begin{frame}{}
  \begin{exampleblock}{$\forallxintro$ 推理规则的应用}
    \[
      \forall x(Q(x)\to R(x)),\ \exists x(P(x)\land Q(x)) \vdash
        \exists x(P(x)\land R(x))
    \]
  \end{exampleblock}

  \begin{logicproof}{2}
    \uncover<2->{\forall x\; (Q(x) \to R(x))} & \uncover<2->{premise} \\
    \uncover<2->{\exists x\; (P(x) \land Q(x))} & \uncover<2->{premise} \\
    \uncover<4->{
    \begin{subproof}
      \uncover<5->{x_0 \;(\text{\blue{fresh}})\; P(x_0) \land Q(x_0)} & \uncover<5->{assumption} \\
      \uncover<10->{Q(x_0)\to R(x_0)} & \uncover<10->{\forallxelim 1} \\
      \uncover<11->{Q(x_0)} & \uncover<11->{\landelimright 3} \\
      \uncover<9->{R(x_0)} & \uncover<12->{\toelim 4, 5} \\
      \uncover<8->{P(x_0)} & \uncover<13->{\landelimleft 3} \\
      \uncover<7->{P(x_0)\land R(x_0)} & \uncover<14->{\landintro 7, 6} \\
      \uncover<6->{\exists x\; (P(x)\land R(x))} & \uncover<15->{\existsxintro 8}
    \end{subproof}}
    \uncover<3->{\exists x\; (P(x)\land R(x))} & \uncover<16->{\existsxelim 2, 3--9}
  \end{logicproof}
\end{frame}
%%%%%%%%%%%%%%%%%%%%

%%%%%%%%%%%%%%%%%%%%
\begin{frame}{}
  \begin{exampleblock}{$\forallxintro$ 推理规则的应用}
    \[
      \Big\{\exists x\; P(x), \forall x \forall y\; (P(x) \to Q(y))\Big\}
        \vdash \forall y\; Q(y)
    \]
  \end{exampleblock}

  \pause
  \begin{logicproof}{2}
    \uncover<2->{\exists x\;P(x)} & \uncover<2->{premise} \\
    \uncover<2->{\forall x\forall y\; (P(x) \to Q(y))} & \uncover<2->{premise} \\
    \uncover<4->{
    \begin{subproof}
      \uncover<5->{y_0} & \uncover<5->{\blue{fresh}} \\
      \uncover<7->{
      \begin{subproof}
        \uncover<8->{x_0\; (\text{\blue{fresh}})} & \uncover<8->{assumption} \\
        \uncover<10->{\forall y(P(x_0)\to Q(y))} & \uncover<10->{\forallxelim 2} \\
        \uncover<11->{P(x_0)\to Q(y_0)} & \uncover<11->{\forallyelim 5} \\
        \uncover<9->{Q(y_0)} & \uncover<12->{\toelim 6, 4}
      \end{subproof}}
      \uncover<6->{Q(y_0)} & \uncover<13->{\existsxelim 1, 4--7}
    \end{subproof}}
    \uncover<3->{\forall y Q(y)} & \uncover<14->{\forallyintro 3--8}
  \end{logicproof}
\end{frame}
%%%%%%%%%%%%%%%%%%%%

%%%%%%%%%%%%%%%%%%%%
\begin{frame}{}
  \begin{exampleblock}{$\existselim$ 推理规则的 (\red{\bf 错误}) 应用}
    \[
      \exists x\; P(x),\ \forall x\; (P(x) \to Q(x)) \vdash \forall y\; Q(y)
    \]
  \end{exampleblock}

  \begin{logicproof}{2}
    \uncover<2->{\exists x P(x)} & \uncover<2->{premise} \\
    \uncover<2->{\forall x(P(x)\to Q(x))} & \uncover<2->{premise} \\
    \uncover<4->{
    \begin{subproof}
      \uncover<5->{x_0} & \uncover<5->{fresh} \\
      \uncover<7->{
      \begin{subproof}
        \uncover<8->{x_{0}\; (\text{\blue{fresh}}),\; P(x_0)} & \uncover<8->{assumption} \\
        \uncover<10->{P(x_0)\to Q(x_0)} & \uncover<10->{\forallxelim 2} \\
        \uncover<9->{Q(x_0)} & \uncover<11->{\toelim 5, 4}
      \end{subproof}}
      \uncover<6->{Q(x_0)} & \uncover<12->{\existsxelim 1, 4--6}
    \end{subproof}}
    \uncover<3->{\forall y\ Q(y)} & \uncover<13->{\forallyintro 3--7}
  \end{logicproof}
\end{frame}
%%%%%%%%%%%%%%%%%%%%

%%%%%%%%%%%%%%%%%%%%
\begin{frame}{}
  \begin{exampleblock}{}
    \[
      \lnot \forall x\; P(x) \vdash \exists x\; \lnot P(x)
    \]
  \end{exampleblock}

  \begin{logicproof}{3}
    \uncover<2->{\lnot \forall x\; P(x)} & \uncover<2->{premise} \\
    \uncover<4->{
    \begin{subproof}
      \uncover<5->{\lnot \exists x\; \lnot P(x)} & \uncover<5->{assumption} \\
      \uncover<8->{
      \begin{subproof}
        \uncover<9->{x_0} & \uncover<9->{fresh} \\
        \uncover<11->{
        \begin{subproof}
          \uncover<12->{\lnot P(x_0)} & \uncover<12->{assumption} \\
          \uncover<14->{\exists x\; \lnot P(x)} & \uncover<14->{\existsxintro 4} \\
          \uncover<13->{\bot} & \uncover<15->{\lnotelim 5, 2}
        \end{subproof}}
      \uncover<10->{P(x_0)} & \uncover<16->{PBC 4-6}
      \end{subproof}}
      \uncover<7->{\forall x\; P(x)} & \uncover<17->{\forallxintro 3--7} \\
      \uncover<6->{\bot} & \uncover<18->{\lnotelim 8, 1}
    \end{subproof}}
    \uncover<3->{\exists x\; \lnot P(x)} & \uncover<19->{PBC 2--9}
  \end{logicproof}
\end{frame}
%%%%%%%%%%%%%%%%%%%%

%%%%%%%%%%%%%%%%%%%%
\begin{frame}{}
  \begin{exampleblock}{}
    \[
      \lnot \forall x\; \phi \vdash \exists x\; \lnot \phi
    \]
  \end{exampleblock}

  \pause
  \begin{logicproof}{3}
    \uncover<2->{\lnot \forall x\; \phi} & \uncover<2->{premise} \\
    \uncover<4->{
    \begin{subproof}
      \uncover<5->{\lnot \exists x\; \lnot \phi} & \uncover<5->{assumption} \\
      \uncover<8->{
      \begin{subproof}
        \uncover<9->{x_0} & \uncover<9->{fresh} \\
        \uncover<11->{
        \begin{subproof}
          \uncover<12->{\lnot \phi[x_0/x]} & \uncover<12->{assumption} \\
          \uncover<14->{\exists x\; \lnot \phi} & \uncover<14->{\existsxintro 4} \\
          \uncover<13->{\bot} & \uncover<15->{\lnotelim 5, 2}
        \end{subproof}}
      \uncover<10->{\phi[x_0/x]} & \uncover<16->{PBC 4-6}
      \end{subproof}}
      \uncover<7->{\forall x\; \phi} & \uncover<17->{\forallxintro 3--7} \\
      \uncover<6->{\bot} & \uncover<18->{\lnotelim 8, 1}
    \end{subproof}}
    \uncover<3->{\exists x\; \lnot \phi} & \uncover<19->{PBC 2--9}
  \end{logicproof}
\end{frame}
%%%%%%%%%%%%%%%%%%%%

%%%%%%%%%%%%%%%%%%%%
\begin{frame}{}
  \begin{exampleblock}{}
    \[
      \exists x\; \lnot \phi \vdash \lnot \forall x\; \phi
    \]
  \end{exampleblock}

  \begin{logicproof}{2}
    \uncover<2->{\exists x\; \lnot \phi} & \uncover<2->{assumption} \\
    \uncover<4->{
    \begin{subproof}
      \uncover<5->{\forall x\; \phi} & \uncover<5->{assumption} \\
      \uncover<7->{
      \begin{subproof}
        \uncover<8->{x_0} & \uncover<8->{fresh} \\
        \uncover<8->{\lnot \phi[x_0/x]} & \uncover<8->{assumption} \\
        \uncover<10->{\phi[x_0/x]} & \uncover<10->{\forallxelim 2} \\
        \uncover<9->{\bot} & \uncover<11->{\lnotelim 5, 4}
      \end{subproof}}
      \uncover<6->{\bot} & \uncover<12->{\existsxelim 1, 3--6}
    \end{subproof}}
    \uncover<3->{\neg \forall x\; \phi} & \uncover<13->{\lnotintro 2--7}
  \end{logicproof}
\end{frame}
%%%%%%%%%%%%%%%%%%%%

%%%%%%%%%%%%%%%%%%%%
\begin{frame}{}
  \begin{exampleblock}{}
    \[
      \exists x\;(\phi\lor\psi) \vdash
        (\exists x\;\phi)\lor(\exists x\;\psi)
    \]
  \end{exampleblock}

  \begin{logicproof}{3}
    \uncover<2->{\exists x\; (\phi \lor \psi)} & \uncover<2->{premise} \\
    \uncover<4->{
    \begin{subproof}
      \uncover<5->{x_0} & \uncover<5->{fresh} \\
      \uncover<5->{(\phi \lor \psi)[x_0/x]} & \uncover<5->{assumption} \\
      \uncover<7->{\phi[x_0/x] \lor \psi[x_0/x]} & \uncover<7->{identical\ to\ 3} \\
      \uncover<8->{
      \begin{subproof}
        \uncover<9->{\phi[x_0/x]} & \uncover<9->{assumption} \\
        \uncover<11->{\exists x\; \phi} & \uncover<11->{\existsxintro 5} \\
        \uncover<10->{(\exists x\;\phi)\lor(\exists x\;\psi)} & \uncover<12->{\lorintroleft 6}
      \end{subproof}}
      \uncover<13->{
      \begin{subproof}
        \uncover<14->{\psi[x_0/x]} & \uncover<14->{assumption} \\
        \uncover<16->{\exists x\; \psi} & \uncover<16->{\existsxintro 8} \\
        \uncover<15->{(\exists x\;\phi)\lor(\exists x\;\psi)} & \uncover<17->{\lorintroright 9}
      \end{subproof}}
      \uncover<6->{(\exists x\;\phi)\lor(\exists x\;\psi)} & \uncover<18->{\lorelim 4, 5--7, 8--10}
    \end{subproof}}
    \uncover<3->{(\exists x\;\phi)\lor(\exists x\;\psi)} & \uncover<19->{\existsxelim 1, 2--11}
  \end{logicproof}
\end{frame}
%%%%%%%%%%%%%%%%%%%%

%%%%%%%%%%%%%%%%%%%%
\begin{frame}{}
  \begin{exampleblock}{\teal{留作练习了}}
    \[
      (\exists x\;\phi)\lor(\exists x\;\psi) \vdash
        \exists x\;(\phi\lor\psi)
    \]

    \pause
    \[
      \exists x\;\exists y\; \phi \vdash
        \exists y\;\exists x\; \phi
    \]
  \end{exampleblock}
\end{frame}
%%%%%%%%%%%%%%%%%%%%

%%%%%%%%%%%%%%%%%%%%
\begin{frame}{\texorpdfstring{$\forall, \exists$}{forall, exists}-推理规则的应用}
  \begin{exampleblock}{请使用一阶谓词逻辑公式表达下列命题与推理}
    给定如下``前提'', 请判断``结论''是否有效, 并说明理由。\\[10pt]
    {\bf 前提:}
    \begin{enumerate}[(1)]
      \item 每个人或者喜欢美剧, 或者喜欢韩剧 (可以同时喜欢二者);
      \item 任何人如果他喜欢抗日神剧, 他就不喜欢美剧;
      \item 有的人不喜欢韩剧。
    \end{enumerate}

    \vspace{0.20cm}
    {\bf 结论:} 有的人不喜欢抗日神剧 ({\it 幸亏如此})。
  \end{exampleblock}

  \pause
  \vspace{0.50cm}
  \[
    \ndnoname{\forall x.\; A(x) \lor K(x) \qquad
      \forall x.\; J(x) \to \lnot A(x) \qquad
      \exists x.\; \lnot K(x)}{\exists x.\; \lnot J(x)}
  \]
\end{frame}

% %%%%%%%%%%%%%%%%%%%%
% \begin{frame}{}
%   \setcounter{equation}{0}
%   \begin{align}
%     \forall x.\; A(x) \lor K(x) & \quad (\text{前提})
%       \label{eq:1} \\[6pt]
%     \forall x.\; J(x) \to \lnot A(x) & \quad (\text{前提})
%       \label{eq:2} \\[6pt]
%     \exists x.\; \lnot K(x) & \quad (\text{前提})
%       \label{eq:3} \\[6pt]
%     \uncover<2->{\red{[x_{0}] \quad [\lnot K(x_{0})]} & \quad (\text{引入变量与假设})
%       \label{eq:4} \\[6pt]}
%     \uncover<3->{A(x_{0}) \lor K(x_{0}) & \quad (\forall\text{-elim}, (\ref{eq:1}), (\ref{eq:4}))
%       \label{eq:5} \\[6pt]}
%     \uncover<4->{A(x_{0}) & \quad ((\ref{eq:4}), (\ref{eq:5}))
%       \label{eq:6} \\[6pt]}
%     \uncover<5->{J(x_{0}) \to \lnot A(x_{0}) & \quad (\forall\text{-elim}, (\ref{eq:2}), (\ref{eq:4}))
%       \label{eq:7} \\[6pt]}
%     \uncover<6->{\lnot J(x_{0}) & \quad ((\ref{eq:6}), (\ref{eq:7}))
%       \label{eq:8} \\[6pt]}
%     \uncover<7->{\red{\exists x.\; \lnot J(x)} & \quad (\exists\text{-intro}, (\ref{eq:8}))
%       \label{eq:9} \\[6pt]}
%     \uncover<8->{\exists x.\; \lnot J(x) & \quad (\exists\text{-elim}, (\ref{eq:3})-(\ref{eq:8}))
%       \label{eq:10}}
%   \end{align}
% \end{frame}
% %%%%%%%%%%%%%%%%%%%%

% %%%%%%%%%%%%%%%%%%%%
% \begin{frame}{}
%   \setcounter{equation}{0}
%   \begin{align}
%     \forall x.\; A(x) \lor K(x) \label{eq:1} \\[3pt]
%     \forall x.\; J(x) \to \lnot A(x) \label{eq:2} \\[3pt]
%     \exists x.\; \lnot K(x) \label{eq:3}
%   \end{align}

%   \pause
%   \begin{center}
%     根据 (\ref{eq:3}), 不妨设 $\lnot K(x)$ 对 $x_{0}$ 成立:
%     \begin{align}
%       \lnot K(x_{0}) \label{eq:4}
%     \end{align}
%     \begin{columns}
%       \column{0.45\textwidth}
%         \pause
%         根据 (\ref{eq:1}), 有
%         \begin{align}
%           A(x_{0}) \lor K(x_{0}) \label{eq:5}
%         \end{align}
%         \pause
%         根据 (\ref{eq:4}) 与 (\ref{eq:5}), 有
%         \begin{align}
%           A(x_{0}) \label{eq:6}
%         \end{align}
%       \column{0.45\textwidth}
%         \pause
%         根据 (\ref{eq:2}), 有
%         \begin{align}
%           J(x_{0}) \to \lnot A(x_{0}) \label{eq:7}
%         \end{align}
%         \pause
%         根据 (\ref{eq:6}) 与 (\ref{eq:7}), 有
%         \begin{align}
%           \lnot J(x_{0}) \label{eq:8}
%         \end{align}
%         \pause
%         因此, $\exists x.\; \lnot J(x)$
%     \end{columns}
%   \end{center}
% \end{frame}
% %%%%%%%%%%%%%%%%%%%%